%% notes-tex/database/database-expansion-hosts/sections/pipeline.tex
%% [[experiments.database.expansion-hosts]]
%% https://github.com/Mjvolk3/torchcell/tree/main/notes-tex/database/database-expansion-hosts/sections/pipeline.tex
%% SOURCE: hand-authored prose. The mirror figures are measured by
%% experiments/database/scripts/audit_provenance_mirror.py and printed in tables/mirror.tex; the
%% quote, gap and schema-field counts come from
%% /scratch/projects/torchcell-scratch/tmp/deep-research/research_notes/Pk and next dataset wave triage/mirror_storage_audit.md,
%% and the API constraints and cost arithmetic from
%% .../research_notes/Pk and next dataset wave triage/artifact_to_ontology_pipeline.md.

\section{From artifact to ontology\secstatus{todo}}\label{sec:pipeline}

%% GENERATED FILE -- do not hand-edit.
%% SOURCE: experiments/database/scripts/audit_provenance_mirror.py
\begin{table}[H]\centering
\footnotesize
\caption[]{What the two provenance mirrors actually store, measured from the trees on disk and from every \texttt{manifest.json} in them, validated through \texttt{torchcell.literature.manifest.Manifest}. \emph{Library} is \texttt{\$DATA\_ROOT/torchcell-library/} (paper PDFs, OCR markdown, SI) and \emph{Raw} is \texttt{\$DATA\_ROOT/torchcell-raw/} (the raw dataset files a loader consumed). Each cell is the share of the population carrying the property, with the count and the population size. Rows are ordered by consequence, worst first: the first block is what a rebuild cannot reproduce, the second what a drift check cannot see, the third the manifest metadata, the fourth shape drift, the last the size of each tree. \emph{Files in no manifest} is the integrity hole: bytes no sha256 pins and no recorded command re-fetches. Byte totals sum file sizes and exclude directory entries, so they read a little under \texttt{du}. \texttt{provenance\_complete} counts manifests asserting a complete chain, so a false value counts as unfilled.}
\label{tab:mirror}
\begin{tabular}{@{}L{74mm} L{48mm} L{48mm}@{}}
\toprule
Property & Library & Raw \\
\midrule
File records with a retrieval record & 2.4\% (599/24,950) & 75.6\% (472/624) \\
Keys where every file record has one & 0.0\% (0/717) & 78.1\% (75/96) \\
Retrieval source ever re-checked (\texttt{last\_check}) & 0.0\% (0/599) & 3.6\% (17/472) \\
File records with a processing record & 1.1\% (278/24,950) & 5.1\% (32/624) \\
\texttt{paper.md} with reproducible MinerU OCR & 15.7\% (112/715) & 0.0\% (0/1) \\
Files on disk named in a manifest & 97.9\% (24,950/25,483) & 92.2\% (624/677) \\
Manifest \texttt{si\_expected} (SI completeness checklist) & 0.0\% (0/717) & 92.7\% (89/96) \\
Manifest \texttt{si\_data\_sources} & 10.5\% (75/717) & 96.9\% (93/96) \\
Manifest \texttt{zotero\_item\_key} & 99.4\% (713/717) & 7.3\% (7/96) \\
Manifest \texttt{collections} & 98.7\% (708/717) & 0.0\% (0/96) \\
Manifest \texttt{doi} & 90.0\% (645/717) & 99.0\% (95/96) \\
Manifest \texttt{provenance\_complete} true & 99.9\% (716/717) & 97.9\% (94/96) \\
File record \texttt{original\_filename} & 2.4\% (590/24,950) & 28.4\% (177/624) \\
File record \texttt{zotero\_md5} & 3.1\% (766/24,950) & 0.0\% (0/624) \\
File record \texttt{source} & 9.2\% (2,291/24,950) & 100.0\% (624/624) \\
File record \texttt{sha256} & 100.0\% (24,950/24,950) & 100.0\% (624/624) \\
Retrieval \texttt{retriever} resolves in \texttt{RETRIEVERS} & 100.0\% (599/599) & 87.5\% (413/472) \\
Retrieval \texttt{source\_url} & 100.0\% (599/599) & 99.8\% (471/472) \\
\midrule
Files in no manifest & 533 files, 0.79 GB (4.9\% of tree) & 53 files, 16.84 GB (48.2\% of tree) \\
Manifest records whose file is gone & 0 & 0 \\
Unresolvable \texttt{retriever} strings & 0 distinct, 0 records & 5 distinct, 59 records \\
File records in the pre-nesting retrieval shape & 2 records in 1 key & 0 records in 0 keys \\
Manifests with fields outside the model & 1 & 1 \\
File roles outside the documented vocabulary & 1 distinct, 2 records & 1 distinct, 1 record \\
Retrieval methods in use & 4 & 9 \\
\midrule
Bytes on disk & 16.10 GB & 34.92 GB \\
Citation-key directories & 717 & 99 \\
Keys carrying a manifest & 100.0\% (717/717) & 97.0\% (96/99) \\
File records across all manifests & 24,950 & 624 \\
Files on disk in key directories & 25,483 & 677 \\
\bottomrule
\end{tabular}
\end{table}


Table~\ref{tab:mirror} measures what the two provenance mirrors store, and the measurement
decides the order of work. In the library mirror 2.4 percent of 24,950 file records carry a
retrieval record, 0 of 717 keys carry one on \texttt{paper.pdf}, whose entire provenance is the
string \texttt{zotero:attachment:<key>}, and 15.7 percent of 715 \texttt{paper.md} files carry
the parser version, arguments, and DPI that produced them. So for 84 percent of the corpus a
re-parse would move every quote anchor with no recorded recipe to reproduce the text that was
quoted. In the raw mirror 48.2 percent of the bytes, 16.84 GB across 19 keys, sit in no
manifest at all, including the 10.2 GB fitness database of Sec.~\ref{sec:pseudo} and all 6.6 GB
of one 2025 single-cell key, and three keys carry no manifest. Nothing re-verifies either tree:
\texttt{last\_check} is set on none of the 599 library retrieval records and on 3.6 percent of
the 472 raw ones, 17 of 1,071 in all, and while \texttt{verify\_artifact} and
\texttt{check\_source} exist and are tested, no
cron, slurm, or CI job calls either, the cron list being five jobs of which none hashes or
re-fetches anything. \texttt{si\_expected}, the per-key completeness checklist, is 0 of 717 in
the library, so no automated check can report that a supplementary item is missing rather than
absent.

What survived into the built stores is the other half of the measurement. Across all 180 built
stores the entire quote corpus is 890 distinct triples of citation key, source sha256, and
verbatim quote, over 95 citation keys and 139 pinned source hashes, against 334,083 typed
\texttt{ProvenanceGap} instances. Absence is recorded at roughly 375 times the rate of
evidence, which is correct policy under the typed-absence rule and is also the measure of how
much of the ontology is unsourced. The reverse edge is thinner still: 1 of 133
\texttt{Publication} constructions sets \texttt{identifier} with a sha256 and the other 132
identify their paper by a DOI or PMID string, \texttt{ArtifactRef} appears in 5 of 123
registered loader modules and reaches 6 of 181 built stores, and \texttt{build\_manifest.json}
carries no mirror reference at all. A served record therefore cannot be resolved to the exact
bytes it was derived from, and a mirror artifact cannot be queried for the records it produced.

\subsection{The nine fields that may carry a sourced value\secstatus{todo}}\label{sec:pipeline-slot}

The binding constraint is narrower and sharper than any of those figures. The schema check
\texttt{sourced\_value\_fields()}, run live, reports exactly nine schema fields that can carry
a \texttt{SourcedValue}, and each of the nine is the \texttt{provenance} field of one class:
\texttt{Media}, \texttt{MediaComponent}, \texttt{PreCulture}, \texttt{CultureFormat},
\texttt{StrainBackground}, \texttt{EukaryoticStrainBackground}, \texttt{BackgroundAllele},
\texttt{IntegratedCassette}, and \texttt{EssentialityStatement}. All nine are media, strain
background, culture format, or essentiality. Nothing on the phenotype value, the perturbation,
or the experiment carries a quote.

That is exactly the output an extraction pipeline produces. A sample count, an uncertainty type,
a unit, and a threshold are the four values the dataset-addition rule requires to be sourced
from the paper or supplement with a verbatim quote, a sha256, and a section or line, and all
four remain loader literals whose justification never leaves the source file. The Kuzmin
combined-mutant sample count is the worked case and it is still a bare literal, with the
back-solve that produced it documented only in a Dendron note, while \texttt{StatDerivation},
the model built for exactly that resolution, is used by 2 of 123 loaders. A pipeline whose
whole value is producing those sourced values has nowhere typed to write them, so the first
thing to build is the slot. A verbatim quote also reaches the LMDB experiment record in only 27
of 181 stores, and for the other 154 it lives in the sibling
\file{preprocess/experiment_reference_index.json}, which is not the served artifact.

\subsection{The staged design\secstatus{todo}}\label{sec:pipeline-stages}

The design that follows from that audit does the unglamorous work first and puts the model only
where its output can be checked. Eight stages, two of them with no model in them at all:

\begin{itemize}[leftmargin=1.2em,itemsep=1pt,topsep=2pt]
  \item \textbf{0. Integrity floor.} No model. Bring the 16.84 GB into the manifest, replace the
    hand-rolled sidecars that no program reads, fix the seven retrieval records that cannot be
    dispatched because their recorded retriever is a shell string, a prose sentence, a
    loader-private function, or a class method, and install the missing re-verify job. A quote
    written against unhashed bytes cannot be checked against anything.
  \item \textbf{1. Parse and pin.} Record the parser name, version, arguments, DPI, and output
    sha256 per artifact. The convention already exists and 112 of 715 keys use it. The parser is
    a replaceable pinned stage rather than a choice, because the benchmark suites disagree with
    each other and the headline figures are vendor-run.
  \item \textbf{2. Grounded read.} A citations-enabled read of the mirrored markdown as a
    \texttt{text/plain} document, emitting claims with character spans into the mirrored bytes.
  \item \textbf{3. Schema-enforced structuring.} A second call that carries the schema, with the
    pydantic surface behind a cache breakpoint.
  \item \textbf{4. Ontology mapping.} A cheap lexical mapper first for recall, then a model
    adjudicating its candidates against a retrieved ontology slice rather than a dumped one.
  \item \textbf{5. Deterministic gate.} No model. Re-hash each cited span against the mirrored
    bytes, and detect truncation, refusal, cache misconfiguration, and ungroundedness from the
    response fields rather than from a judgment. This is \texttt{audit\_sourced\_value}, which
    already exists, is called from 31 of 123 loader modules, and is skipped in CI by a sentinel
    data root.
  \item \textbf{6. Adjudication.} The expensive model, on the disagreements only, so its spend
    stays proportional to the failure rate.
  \item \textbf{7. Human review.} Unavoidable, and quantified below.
\end{itemize}

Three API constraints determine where those stage boundaries fall. Citations and structured
outputs are mutually exclusive in one request and return a 400 together, so grounding and
schema enforcement are two calls and not one. Native PDF input is vision plus per-page
extracted text rather than an OCR service, and it cannot cite images in a PDF at all, so a
figure-borne number cannot be grounded that way. Plain text chunks into sentences and yields
character spans into bytes the mirror pins, where a PDF yields page spans into bytes the
publisher controls. And the schema handed to the API is not the pydantic model verbatim,
because recursive schemas, \texttt{oneOf}, \texttt{not}, numerical constraints, and
string-length constraints each return a 400, so the SDK helpers strip the unsupported
constraints and re-validate locally. A supplementary spreadsheet is rejected outright in a
document block, so it becomes CSV text, which keeps the citation span, or an upload into code
execution, which loses it.

\subsection{The cost model, and its assumptions\secstatus{todo}}\label{sec:pipeline-cost}

The figures below are arithmetic over the documented prices, and the assumptions are stated
because the result moves with them:

\begin{itemize}[leftmargin=1.2em,itemsep=1pt,topsep=2pt]
  \item 300 papers, each a 12-page main PDF plus a 20-page supplement plus about 30,000 tokens
    of converted supplementary tables.
  \item 3,500 tokens per page, the midpoint of the documented 1,500 to 3,000 text tokens plus
    up to 1,568 visual tokens, which is about 150,000 input tokens of artifact per paper.
  \item A 20,000-token schema prefix, cached.
  \item Outputs of 8,000, 4,000, 2,000, and 15,000 tokens for the four model stages.
  \item The adjudication stage firing on 20 percent of papers, and a 30 percent retry
    overhead on everything.
\end{itemize}

On those assumptions 300 papers cost about \$120 to \$175 with a mid-tier worker and a
top-tier adjudicator run in batch, and the planning figure is \$190 to \$250. Running every
stage on the top-tier model synchronously at high effort costs about \$810 to \$1,000 for the
same work, roughly four to five times as much. Input tokens on the artifact are 70 to 80
percent of the bill in every variant, so one design detail is load-bearing. Caching the
20,000-token schema prefix saves about \$0.02 per paper, while caching the 150,000-token
document between stages saves about \$0.14, about a third of the total. The documentation states
that an ephemeral cache entry written during batch processing would likely expire before the
follow-up request runs, so a batch-based multi-pass design cannot count on cross-batch hits and
either puts all passes for one paper in a single request or budgets the higher figure.
Batch-first is a throughput choice before it is a cost one, since the batch API is documented
as exempt from the ramp limits that return a 429 on a burst even while under the token limits,
and cache reads do not count toward the input-token limit either.

The accuracy regimes are the number worth planning against. Published pipelines of this shape
land at 0.94 to 0.98 F1, with one entity and relation extraction at 0.96 and 0.94 and a
protein-degradation workflow at record-level F1 0.98 with 92 percent and 82.5 percent of
recovered records validated on expert review. A bare model on strict whole-record extraction
lands at 55 to 60 percent, measured at 57 percent and 56 percent for two current general
models on a nanocrystal-synthesis benchmark. At 300 papers and about 20 fields each, the good
regime still implies 120 to 360 wrong fields, so the pipeline's value is in making a wrong
field findable rather than in removing the review queue. The gap between the two regimes is the
engineering and not the model.

Two targeting decisions follow from the audit rather than from the API. The pipeline points at
the raw tree and this wave's keys rather than at the library mirror, because 124 of 717 library
keys are named anywhere in \file{torchcell/} source and only 95 are quoted by a built store,
so roughly 554 keys and about 12 GB are reading material served through the literature API
rather than substrate the ontology depends on, while 96 of 99 raw keys are named in code. And
the cheapest genuinely useful first output is not new records: \texttt{ProvenanceGap.resolve\_with},
the field that names the artifact which would close a deferred gap, is populated on 0.8 percent
of gaps against 4,933 gaps typed \texttt{deferred\_pending\_source\_review}. Filling it needs
no schema change, has a deterministic check, and turns an existing 4,933-item worklist into an
actionable one.

\subsection{What this wave needs changed\secstatus{todo}}\label{sec:pipeline-schema}

Table~\ref{tab:schema} lists every change in the served vocabulary and the candidate machinery
that the rows of this document require, with the policy each one carries. Two of them gate
every row rather than one: a host token the gate can resolve, without which
\texttt{hosts\_named()} raises on any \org{P. kudriavzevii} row so that G2 cannot even be
measured, and an \texttt{Organism} member, without which no loader can declare the host and the
adapter map can hold nothing for it. The most expensive entry is the slot of
Sec.~\ref{sec:pipeline-slot}: extending \texttt{SourcedValue} to the phenotype value and the
perturbation changes classes that many served closures hold, so its rebuild scope reads
\texttt{full\_kg\_rebuild} and its record carries the \texttt{schema-breaking} label. Every
other entry on the table is additive, meaning a new leaf or a new member that changes no record
any served dataset already holds.

%% GENERATED FILE -- do not hand-edit.
%% SOURCE: experiments/database/scripts/build_hosts_candidate_datasets_table.py
\begingroup
\footnotesize
\begin{longtable}{@{}L{50mm} L{48mm} L{74mm}@{}}
\caption[]{Every schema or candidate-machinery change this wave requires, with the
symbol it lands on and the policy it triggers. \emph{Symbol} is the exact module path
and name, so the change is addressable rather than described. A change marked additive
adds a class or a vocabulary member that only its own closure holds, so it enters the
served graph by incremental admission; a change to a class other closures hold is a full
rebuild of every dataset whose closure holds it. The first two are the critical path:
without them no loader can declare a host and the candidate gate cannot even be
measured.}
\label{tab:schema}\\
\toprule
Symbol & What changes & Policy it triggers \\
\midrule
\endfirsthead
\toprule
Symbol & What changes & Policy it triggers \\
\midrule
\endhead
\bottomrule
\endfoot
\file{torchcell/datamodels/schema.py::Genotype (perturbation identity)} & Admit a perturbation whose identity is NOT a gene. Eleven of the twenty largest yeast rows in this fill, and most of the E. coli regulatory-DNA rows, carry no gene identifier anywhere in their bulk data: the record key is a lineage barcode, a raw measured sequence, or a residue substitution string. A barcoded lineage is the hardest case, because the perturbation is UNKNOWN by construction and is identified only by its tag, which is a typed absence on the genotype rather than a missing field & Served-vocabulary change. The honest route is a typed ProvenanceGap on the gene-identity field rather than a forced mapping to a gene, which would invent a locus the release never names. Schema-change record and impact report; additive if it lands as a new perturbation sibling, breaking if it loosens a required field on the served Genotype Not additive. \\
\addlinespace[5pt]
\file{torchcell/datamodels/schema.py::SourceType} & Add a patent member. The enum holds journal\_article, dissertation, preliminary\_report and in\_house, and the identity rule for every non-journal member is title plus the deposited document's mirror-relative path and sha256. A patent fits that rule exactly and has no DOI, so the three ranked 3-HP patents can be deposited and hashed but not yet declared. Without the member the only honest alternatives are to drop eleven strains with titer and yield, or to file a patent as a journal article, which would be false & Served-vocabulary change, but additive: no served record carries the new member, so the closure of every built dataset is unchanged. Schema-change record and impact report; no schema-breaking label and no rebuild, since nothing served reserializes Additive. \\
\addlinespace[5pt]
\file{torchcell/data/experiment_dataset.py::Organism} & Add a pkudriavzevii member. The enum holds exactly scerevisiae, ecoli and pputida today, and ALL\_ORGANISMS is frozenset(Organism), so no loader can declare the host and dataset\_adapter\_map can hold nothing for it & Served-vocabulary change: schema-change record, impact report, the served datasets named in the PR, and the schema-breaking label if a served closure changes. The genome side is already done, which is what makes this the whole cost: registry.py pins PKUDRIAVZEVII\_SD108 = pkudriavzevii\_SD108\_v1 and gates.py already lists it in READABLE\_ASSEMBLY\_SETS with the comment that no host token names it yet. Additive. \\
\addlinespace[5pt]
\file{torchcell/candidates/gates.py::HOSTS, _ORGANISM_OF_TABLE_ROW} & Add HostTokens for SD108 and CBS573, and admit 'P. kudriavzevii' to the set of organisms a table row may name, which holds E. coli, P. putida and S. cerevisiae today & Gate-code change, not a schema change. Until it lands hosts\_named() raises ValueError on any P. kudriavzevii row, so G2 cannot even be measured and no verdict can be written for one. The eight pseudomonad hosts need the same treatment, one token per reference strain. Additive. \\
\addlinespace[5pt]
\file{torchcell/data/experiment_dataset.py::Organism (pseudomonads)} & Nine or ten further members for the Fitness Browser organisms and P. aeruginosa, plus the matching reference genomes in the genomes tier & The same served-vocabulary policy, ten times over. Worth staging as one deliberate change rather than ten accidental ones, because the Fitness Browser rows arrive in one artifact and the marginal work per organism after the first is a genome registration. Additive. \\
\addlinespace[5pt]
\file{torchcell/candidates/verdict.py::CandidateTable} & A third literal for this pool. The type is Literal['bacteria', 'yeast'] and it is also the table field of every one of the 52 stored verdicts & Record-model change on a frozen, extra-forbidding model, so it is a policy decision rather than a widening: until it lands, every row in Table 11 reads 'no verdict' and candidate-gate can read this pool's rows but cannot write a verdict for one. gates.py::CandidateRow.table and gates.py::TEXT\_FIELDS carry the same literal and the same keys. Not additive. \\
\addlinespace[5pt]
\file{torchcell/candidates/cli.py::TABLE_SCRIPTS} & Register this script as a third pool, so candidate-gate can load its rows by name the way it loads the bacteria and yeast pools & CLI change only. Done. The gate subcommand's --table choices still admit two values, so passing this pool's name needs the CandidateTable decision above first. Additive. \\
\addlinespace[5pt]
\file{experiments/database/scripts/build_candidate_datasets_table.py::CANDIDATES} & A separate pool rather than more rows in the yeast one. The yeast Candidate has no organism field at all, the bacteria one admits only E. coli and P. putida, and the yeast main() raises on a pool over 250 rows against 169 literals held & Splitting by wave is forced by the guard, not merely tidier. This script is that split, and its own guard is 400 because the pool spans eight hosts. Additive. \\
\addlinespace[5pt]
\file{experiments/database/scripts/build_hosts_candidate_datasets_table.py::SeqBasis} & Host-specific sequence bases. The yeast SeqBasis is eleven values, every one S288C- or yeast-panel-relative, so isolate-WGS transfers to a Pichia or pseudomonad isolate panel and an SD108 designed edit does not & Script-level vocabulary, enforced as a validation error on any unlisted value. Declared explicitly per host so a bacterial row cannot be written against a yeast-shaped basis by accident. Additive. \\
\addlinespace[5pt]
\file{torchcell/datamodels/schema.py (phenotype leaves)} & New leaves for TSS initiation efficiency, antifungal and antibiotic MIC, Biolog respiration per well, plate-reader reporter ratio, and 13C flux & By the placement rule a different technology is its own leaf with an explicit conversion, so a plate-reader GFP ratio and a sort-seq bin average are different leaves rather than one class with a note. Each is additive and admissible on its own and none weakens an existing required field. Additive. \\
\addlinespace[5pt]
\file{torchcell/datamodels/schema.py (perturbation leaves)} & New leaves for a 5' UTR variant, a 3' end variant, a promoter swap upstream of a native gene, a tRNA point substitution, and an amino-acid substitution for the ProteinGym subset & Additive leaves, so each should clear the admission check on its own closure. In the served graph the perturbation subclasses collapse into one Perturbation node whose kind is a property, so a query naming every deletion names the kinds one by one. Additive. \\
\addlinespace[5pt]
\file{torchcell/datamodels/schema.py (GenePerturbation namespaces)} & Gene-identifier namespaces for a JGI gNNNN model id, an INSDC SD108 locus tag, a PA or PA14 number, a Bakta locus tag and a Fitness Browser locusId & The same additive change the bacterial layer already made for b-numbers and PP\_ tags. The genome never gates a record, so a released name the reference does not hold is a report-class verification finding and never a dropped row. Additive. \\
\addlinespace[5pt]
\file{torchcell/verification/sourced.py::SourcedValue} & Extend it to the phenotype value and the perturbation. ontology\_checks.sourced\_value\_fields() reports exactly nine fields that can carry one today, all of them media, strain background, culture format or essentiality, so a quoted n\_samples, uncertainty type, unit or threshold has nowhere typed to be written & The most expensive change here: it moves classes many served closures hold, so the record's rebuild\_scope reads full\_kg\_rebuild and the PR carries the schema-breaking label. It is also the blocking change for an artifact-to-ontology pipeline, whose entire output is that class of value. Not additive. \\
\addlinespace[5pt]
\file{notes/paper.north-star.dataset-triage.md (guide-level exclusion)} & An explicit decision on gRNA-library rows. The standing rule excludes a dataset whose primary label is per-guide activity or cleavage efficiency rather than a per-strain phenotype & The landing-pad screen's cutting and integration efficiency fall under the rule while its fitness and GFP readouts do not, and the PAO1 and KT2440 CRISPRi libraries are guide-level by construction, where the guide sequence is part of the genotype. This needs a decision, not an inherited rule. Additive. \\
\addlinespace[5pt]
\file{experiments/database/scripts/build_candidate_datasets_table.py::EXCLUDED} & Revisit the standing off-species exclusion of the K. phaffii screen and the Y1000+ panel, whose note reads 'Revisit if the generalization axis widens to other hosts' & A cross-host wave is that widening, so the note needs a decision. The iIsor850 exclusion stands on its own merits under not-a-dataset: a reconstruction is not an experiment. Additive. \\
\addlinespace[5pt]
\file{$DATA_ROOT/torchcell-raw/priceMutantPhenotypesThousands2018} & Bring the key into the manifest before extending it. Six of its nine files are unhashed over 10.21 GB, and a parallel non-pydantic convention sits beside it: DEPOSIT.md, feba.db.sha256 and SHA256SUMS.txt, which no program reads & A precondition of the cross-host fitness band rather than a schema change: those rows extend this key, and an ingested fitness value has to carry the release date and sha256 of the db file it came from, because the bulk page itself says the 2018 release's sucrose and D-mannitol data should be disregarded in favor of fresh stock data. Additive. \\
\addlinespace[5pt]
\end{longtable}
\endgroup

